\documentclass[10pt, letterpaper]{article}

% Packages:
\usepackage[
    ignoreheadfoot, % set margins without considering header and footer
    top=0.5in, % separation between body and page edge from the top
    bottom=0.5in, % separation between body and page edge from the bottom
    left=0.5in, % separation between body and page edge from the left
    right=0.5in, % separation between body and page edge from the right
    footskip=1.0 cm, % seperation between body and footer
    % showframe % for debugging
]{geometry} % for adjusting page geometry
\usepackage{titlesec} % for customizing section titles
\usepackage{tabularx} % for making tables with fixed width columns
\usepackage{array} % tabularx requires this
\usepackage[dvipsnames]{xcolor} % for coloring text
\definecolor{primaryColor}{RGB}{0, 79, 144} % define primary color
\usepackage{enumitem} % for customizing lists
\usepackage{fontawesome5} % for using icons
\usepackage{amsmath} % for math
\usepackage{comment}
\usepackage[
    pdftitle={Hector Hernandez - ML Engineer Resume},
    pdfauthor={Hector Hernandez},
    pdfcreator={LaTeX with RenderCV},
    colorlinks=true,
    urlcolor=primaryColor
]{hyperref} % for links, metadata and bookmarks
\usepackage[pscoord]{eso-pic} % for floating text on the page
\usepackage{calc} % for calculating lengths
\usepackage{bookmark} % for bookmarks
\usepackage{lastpage} % for getting the total number of pages
\usepackage{changepage} % for one column entries (adjustwidth environment)
\usepackage{paracol} % for two and three column entries
\usepackage{ifthen} % for conditional statements
\usepackage{needspace} % for avoiding page brake right after the section title
\usepackage{iftex} % check if engine is pdflatex, xetex or luatex
\usepackage[none]{hyphenat} % disable automatic hyphenation

% Ensure that generate pdf is machine readable/ATS parsable:
\ifPDFTeX
    \input{glyphtounicode}
    \pdfgentounicode=1
    \usepackage[utf8]{inputenc}
    \usepackage[T1]{fontenc}
    \IfFileExists{uarial.sty}{
        \usepackage{uarial}
    }{
        \usepackage{helvet}
    }
    \renewcommand{\familydefault}{\sfdefault}
\else
    \usepackage{fontspec}
    \setsansfont{Arial}
    \setmainfont{Arial}
    \renewcommand{\familydefault}{\sfdefault}
\fi



% Some settings:
\AtBeginEnvironment{adjustwidth}{\partopsep0pt} % remove space before adjustwidth environment
% \pagestyle{empty} % no header or footer (commented out to allow custom footer style)
\setcounter{secnumdepth}{0} % no section numbering
\setlength{\parindent}{0pt} % no indentation
\setlength{\topskip}{0pt} % no top skip
\setlength{\columnsep}{0cm} % set column seperation
\hyphenpenalty=10000
\exhyphenpenalty=10000
\makeatletter
\let\ps@customFooterStyle\ps@plain % Copy the plain style to customFooterStyle
\patchcmd{\ps@customFooterStyle}{\thepage}{
    \color{gray}\textit{\small Hector Hernandez - Page \thepage{} of \pageref*{LastPage}}
}{}{} % replace number by desired string
\makeatother
\pagestyle{customFooterStyle}

\titleformat{\section}{\needspace{4\baselineskip}\bfseries\large}{}{0pt}{}[\vspace{1pt}\titlerule]

\titlespacing{\section}{
    % left space:
    -1pt
}{
    % top space:
    0.08 cm
}{
    % bottom space:
    0.12 cm
} % section title spacing

\renewcommand\labelitemi{$\circ$} % custom bullet points
\newenvironment{highlights}{
    \begin{itemize}[
        topsep=0.10 cm,
        parsep=0.10 cm,
        partopsep=0pt,
        itemsep=0pt,
        leftmargin=0.4 cm + 10pt
    ]
}{
    \end{itemize}
} % new environment for highlights

\newenvironment{highlightsforbulletentries}{
    \begin{itemize}[
        topsep=0.10 cm,
        parsep=0.10 cm,
        partopsep=0pt,
        itemsep=0pt,
        leftmargin=10pt
    ]
}{
    \end{itemize}
} % new environment for highlights for bullet entries


\newenvironment{onecolentry}{
    \begin{adjustwidth}{
        0.2 cm + 0.00001 cm
    }{
        0.2 cm + 0.00001 cm
    }
}{
    \end{adjustwidth}
} % new environment for one column entries

\newenvironment{twocolentry}[2][]{
    \onecolentry
    \def\secondColumn{#2}
    \setcolumnwidth{\fill, 4.5 cm}
    \begin{paracol}{2}
}{
    \switchcolumn \raggedleft \secondColumn
    \end{paracol}
    \endonecolentry
} % new environment for two column entries

\newenvironment{header}{
    \setlength{\topsep}{0pt}\par\kern\topsep\centering\linespread{1.5}
}{
    \par\kern\topsep
} % new environment for the header

\newcommand{\placelastupdatedtext}{% \placetextbox{<horizontal pos>}{<vertical pos>}{<stuff>}
  \AddToShipoutPictureFG*{% Add <stuff> to current page foreground
    \put(
        \LenToUnit{\paperwidth-2 cm-0.2 cm+0.05cm},
        \LenToUnit{\paperheight-1.0 cm}
    ){\vtop{{\null}\makebox[0pt][c]{
        %\small\color{gray}\textit{Last updated in September 2024}\hspace{\widthof{Last updated in September 2024}}
    }}}%
  }%
}%

% save the original href command in a new command:
\let\hrefWithoutArrow\href

% new command for external links:
\renewcommand{\href}[2]{\hrefWithoutArrow{#1}{\ifthenelse{\equal{#2}{}}{ }{#2 }\raisebox{.15ex}{\footnotesize \faExternalLink*}}}


\begin{document}
    \newcommand{\AND}{\unskip
        \cleaders\copy\ANDbox\hskip\wd\ANDbox
        \ignorespaces
    }
    \newsavebox\ANDbox
    \sbox\ANDbox{}

    \placelastupdatedtext
    \begin{header}
        \textbf{\fontsize{24 pt}{24 pt}\selectfont Hector Hernandez}

        \vspace{0.3 cm}

        \normalsize
        \mbox{{\color{black}\footnotesize\faMapMarker*}\hspace*{0.13cm}Chandler, AZ 85224}%
        \kern 0.25 cm%
        \AND%
        \kern 0.25 cm%
        \mbox{\hrefWithoutArrow{mailto:hhhector9@gmail.com}{\color{black}{\footnotesize\faEnvelope[regular]}\hspace*{0.13cm}hhhector9@gmail.com}}%
        \kern 0.25 cm%
        \AND%
        \kern 0.25 cm%
        \mbox{\hrefWithoutArrow{tel:+1-480371-5537}{\color{black}{\footnotesize\faPhone*}\hspace*{0.13cm}480-371-5537}}%
        \kern 0.25 cm%
        \AND%
        \kern 0.25 cm%
        \mbox{{\color{black}\footnotesize\faFlag}\hspace*{0.13cm}US Citizen}%
        \kern 0.25 cm%
        \AND%
        \kern 0.25 cm%
        \mbox{\hrefWithoutArrow{https://www.hectorhernandez.dev}{\color{black}{\footnotesize\faLink}\hspace*{0.13cm}hectorhernandez.dev}}%
        \kern 0.25 cm%
        \AND%
        \kern 0.25 cm%
        %\mbox{\hrefWithoutArrow{https://linkedin.com/in/yourusername}{\color{black}{\footnotesize\faLinkedinIn}\hspace*{0.13cm}yourusername}}%
        \kern 0.25 cm%
        \AND%
        \kern 0.25 cm%
        %\mbox{\hrefWithoutArrow{https://github.com/yourusername}{\color{black}{\footnotesize\faGithub}\hspace*{0.13cm}yourusername}}%
    \end{header}

    \vspace{0.3 cm - 0.3 cm}


    \section{Summary}

        \begin{onecolentry}
            ML/Agentic AI Engineer with 10 years in Intel production environments and hands-on experience building LLM-based tools, multi-agent systems, and scalable data platforms. I design reliable agentic systems that combine models, tools, and data into end-to-end multi-step workflows with explicit planning, structured outputs, tool use, failure handling, and monitoring and observability. Currently completing an MS in Computer Science at ASU while applying LLM-assisted coding, evaluation, and latency/cost/accuracy tradeoffs to ship practical AI systems.
        \end{onecolentry}

        \vspace{0.008 cm}

    \section{Education}

            \begin{twocolentry}{
        \textit{Expected December 2026}}
            \textbf{Arizona State University}

            \textit{MS in Computer Science (GPA: 3.88/4.0)}
        \end{twocolentry}

        \vspace{0.10 cm}
        \begin{onecolentry}
            \begin{highlights}
                \item \textbf{Coursework:} Agentic AI, Statistical Machine Learning, Artificial Intelligence, Data Processing at Scale, Advanced Operating Systems, Data Structures and Algorithms, Applied Cryptography, Data Visualization.
            \end{highlights}
        \end{onecolentry}

\vspace{0.08 cm}

        \begin{twocolentry}{
        \textit{Aug 2012 -- May 2015}}
            \textbf{Arizona State University}

            \textit{BS in Chemical Engineering}
        \end{twocolentry}

    \vspace{0.15 cm}

    \section{Experience}

        \begin{twocolentry}{
        \textit{Atlanta, GA}

        \textit{May 2026 -- August 2026}}
            \textbf{Data Engineering Intern}

            \textit{Honeywell}
        \end{twocolentry}

        \vspace{0.10 cm}
        \begin{onecolentry}
            \begin{highlights}
                \item Building Python and SQL data pipelines on Databricks across Azure and GCP data lakes that support analytics and reporting workflows across business units, with emphasis on scalable data models and downstream reliability.
                \item Deploying pipelines and services through CI/CD on Kubernetes, automating build, test, and release to improve reliability and reduce manual intervention.
                \item Developing LLM-based agents for industrial engineering data analysis, combining tool use and structured outputs to automate exploration and surface insights from large operational datasets.
                \item Collaborating with engineering teams to improve data quality, automate manual processing, and ship maintainable backend systems.
            \end{highlights}
        \end{onecolentry}

        \vspace{0.2 cm}

        \begin{twocolentry}{
        \textit{Chandler, AZ}

        \textit{June 2020 -- July 2025}}
            \textbf{Software Application Engineer}

            \textit{Intel}
        \end{twocolentry}

        \vspace{0.10 cm}
        \begin{onecolentry}
            \begin{highlights}
                \item Designed end-to-end architecture for Intel's global EDI manufacturing data platform, integrating Python prototypes, Angular/D3/Highcharts visualizations, and backend services for thousands of engineering users.
                \item Re-architected a legacy Python/Oracle ETL stack into C\#/SQL Server, making cost/performance tradeoffs that improved scalability, maintainability, and long-term operational efficiency.
                \item Optimized database structures and query paths across services, reducing data transfer latency from 30+ seconds to under 3 seconds for production engineering workflows.
                \item Built Airflow, Docker, Python, and SQL Server pipelines with clear failure handling and reduced manual intervention, improving reliability of recurring manufacturing data workloads.
                \item Maintained globally distributed production infrastructure, resolving bugs and enhancements while strengthening monitoring, observability, and availability for manufacturing teams.
                \item Led code reviews and Python standards for the EDI development team, using LLM-assisted coding and structured review habits to improve quality, testability, and defect prevention.
            \end{highlights}
        \end{onecolentry}

        \vspace{0.2 cm}

        \begin{twocolentry}{
        \textit{Chandler, AZ}

        \textit{June 2015 -- May 2020}}
            \textbf{Photolithography Process Engineer}

            \textit{Intel}
        \end{twocolentry}

        \vspace{0.10 cm}
        \begin{onecolentry}
            \begin{highlights}
                \item Sustained and improved Photolithography processes to meet manufacturing demand for throughput, cycle time, and yield.
                \item Improved quality levels by reducing scrap, minimizing yield loss mechanisms, and enhancing process control.
                \item Identified and executed cost-reduction solutions aligned with factory budget goals.
                \item Collaborated with Manufacturing, Maintenance, and Process teams to optimize the overall photolithography area, working extensively with both Nikon and ASML scanners.
                \item Reduced defectivity mechanisms through targeted engineering solutions and tighter process control.
                \item Developed and deployed automation tools, including 1-click reports, SQL queries, JSL scripts, and Python scripts to extract, process, and visualize large datasets from legacy and modern systems.
                \item Trained and mentored multiple new engineers, providing hands-on guidance to help them understand semiconductor equipment and processes, accelerating their ramp-up in the photolithography area.
        \end{highlights}
        \end{onecolentry}

    \section{AI Projects \& Research}

        \begin{twocolentry}{
        \textit{\href{https://github.com/HectorHernandez1/VoiceValet}{GitHub}}}
            \textbf{VoiceValet}
        \end{twocolentry}
        \vspace{0.10 cm}
        \begin{onecolentry}
            \begin{highlights}
                \item Built an agentic phone agent using Claude, Vapi, Deepgram, and Resend to execute outbound calls, integrate tools, handle multi-step tasks, recover from failures, and return transcript summaries by email.
            \end{highlights}
        \end{onecolentry}

        \vspace{0.2 cm}

        \begin{twocolentry}{
        \textit{\href{https://github.com/HectorHernandez1/ConvergeAI}{GitHub}}}
            \textbf{ConvergeAI}
        \end{twocolentry}
        \vspace{0.10 cm}
        \begin{onecolentry}
            \begin{highlights}
                \item Built a multi-model consensus solver that runs GPT and Claude in parallel, compares disagreement, refines answers, caches responses, and tracks cost for latency/cost/accuracy tradeoffs.
            \end{highlights}
        \end{onecolentry}

        \vspace{0.2 cm}

        \begin{twocolentry}{
        \textit{\href{https://github.com/HectorHernandez1/ClaudeMCP}{GitHub}}}
            \textbf{Claude MCP Server Collection}
        \end{twocolentry}
        \vspace{0.10 cm}
        \begin{onecolentry}
            \begin{highlights}
                \item Created 5 MCP servers for stocks, weather, news, Gmail, and PostgreSQL finance data, extending Claude with real-time API access, read-only queries, OAuth actions, and graceful error handling.
            \end{highlights}
        \end{onecolentry}

        \vspace{0.2 cm}

        \begin{twocolentry}{
        \textit{\href{https://github.com/HectorHernandez1/NexusResearch}{GitHub}}}
            \textbf{Nexus Research}
        \end{twocolentry}
        \vspace{0.10 cm}
        \begin{onecolentry}
            \begin{highlights}
                \item Built a multi-agent research pipeline that decomposes questions, retrieves public API context, checks reliability, and emits structured Markdown reports.
            \end{highlights}
        \end{onecolentry}

        \vspace{0.2 cm}

        \begin{twocolentry}{
        \textit{CSE 598 Research}}
            \textbf{$\tau$-bench Research}
        \end{twocolentry}
        \vspace{0.10 cm}
        \begin{onecolentry}
            \begin{highlights}
                \item Evaluated agentic tool use across Qwen3 and Gemma4 families using $\tau$-bench tasks, error taxonomies, $pass^k$ metrics, vLLM, SLURM, and HPC.
            \end{highlights}
        \end{onecolentry}

    \section{Technologies}

        \begin{onecolentry}
            \textbf{Agentic AI \& LLMs:} Claude API, OpenAI API, MCP, LangChain, Ollama, vLLM, ReAct, Function Calling, multi-agent orchestration, tool-use evaluation
        \end{onecolentry}

        \vspace{0.2 cm}

        \begin{onecolentry}
            \textbf{Backend \& Systems:} Python, FastAPI, Docker, Kubernetes, PostgreSQL, SQL Server, Airflow, Databricks, CI/CD, C\#, .NET
        \end{onecolentry}

        \vspace{0.2 cm}

        \begin{onecolentry}
            \textbf{ML \& Evaluation:} scikit-learn, Pandas, NumPy, $\tau$-bench, WildToolBench, $pass^k$ metrics
        \end{onecolentry}

        \vspace{0.2 cm}

        \begin{onecolentry}
            \textbf{Dev \& Infra:} Git, GitHub, Claude Code CLI, SLURM, HPC (A100 / AMD ROCm), VS Code
        \end{onecolentry}

        \vspace{0.2 cm}

        \begin{onecolentry}
            \textbf{Supporting:} JavaScript, TypeScript, React, Angular, D3.js
        \end{onecolentry}

\end{document}
